% Save as: module7-problem-set.tex
% ============================================================================
% Problem Set 7: Development, Labor & Education
% Microeconomic Theory for Experimentalists & Applied Microeconomists
% Built from templates/problem_set_template.tex. Compile with pdflatex.
% ============================================================================
\documentclass[11pt]{article}
\usepackage{amsmath, amssymb, amsthm}
\usepackage[margin=1in]{geometry}
\usepackage{enumitem}

\newtheorem{question}{Question}

\title{Problem Set 7: Development, Labor \& Education\\
\large Microeconomic Theory for Experimentalists \& Applied Microeconomists}
\author{Due: \textbf{[date --- before Lecture 1 of Module 8]}}
\date{}

\begin{document}
\maketitle

\noindent\emph{Ground rules: collaboration on ideas is encouraged; write-ups
are individual. You may use AI tools for parts flagged with
$\langle$AI$\rangle$ and must attach the transcript. Show all calculations.}

\medskip

\noindent\textbf{Weights:} Part A --- 30 points. Part B --- 45 points.
Part C --- 25 points.

\bigskip

% ---------------------------------------------------------------------------
\section*{Part A: Theory --- Guided Calculations (30 points)}

\begin{question}[Dynamic complementarity with numbers --- 15 points]
A child's adult skill depends on early-childhood investment $I_E$ and
later investment $I_L$. Consider two candidate technologies:
\emph{complements}, $S = I_E \times I_L$, and \emph{substitutes},
$S = I_E + I_L$. A family (or a government) has 10 units to allocate.

\begin{enumerate}[label=(\alph*)]
  \item \textbf{(4 pts)} Under each technology, find the allocation of
        the 10 units that maximizes $S$, and the resulting skill level
        (if the maximizer is not unique, say so).
        \emph{Hint: under complements, think symmetry.}
  \item \textbf{(4 pts)} Now suppose early investment is stuck at
        $I_E = 2$ (the early window was missed). Under the complements
        technology, how many units of late investment are needed to
        reach the skill level 25 that a balanced 5/5 allocation would
        have produced? What is the total cost of that remediation path
        (count both the early and the late units), and what is the
        marginal product of a late unit at $I_E = 2$ versus at
        $I_E = 5$?
  \item \textbf{(4 pts)} In one paragraph: why do late-remediation
        programs systematically underperform early programs \emph{if}
        the technology is complementary, and what 2$\times$2 experiment
        identifies whether it actually is? State the estimand in one
        phrase (Module 6 named it).
  \item \textbf{(3 pts)} A skeptic notes that production-function
        parameters estimated from observational data are contaminated
        by \emph{parental responses}: parents may reinforce early
        advantages or compensate for early deficits. Explain in two
        sentences how each response biases the observational estimate
        of complementarity, and why randomization in the 2$\times$2
        experiment of (c) removes this bias (parental responses to the
        programs themselves remain part of the measured effect ---
        Part B(c) returns to this).
\end{enumerate}
\end{question}

\begin{question}[Deworming, externalities, and cost-effectiveness --- 15
points]
A deworming program costs \$0.50 per treated child per year. Evaluation
evidence (stylized Miguel--Kremer magnitudes) shows each treated child
gains 0.075 school-years per year of treatment. In addition, each
treated child generates an infection-reduction externality: 2 nearby
untreated children each gain 0.02 school-years.

\begin{enumerate}[label=(\alph*)]
  \item \textbf{(5 pts)} Compute the cost per additional school-year
        (i) counting only the direct effect and (ii) counting direct
        plus externality effects. By what factor does ignoring the
        externality overstate the cost?
  \item \textbf{(4 pts)} The externality was only visible because
        randomization was at the \emph{school} level. Explain in three
        sentences why child-level randomization within schools would
        have (i) contaminated the control group and (ii) biased the
        estimated treatment effect, stating the direction of the bias.
  \item \textbf{(3 pts)} A ministry ranks programs by cost per
        school-year and funds the top three. Using your numbers,
        explain how the choice of randomization level in the
        \emph{evaluation} can change which programs get \emph{funded}
        --- one concrete sentence about the ranking.
  \item \textbf{(3 pts)} Name one other intervention type (any domain)
        whose evaluation plausibly zeroed out a spillover by
        randomizing at too fine a level, and say in one sentence what
        the coarser design would be.
\end{enumerate}
\end{question}

% ---------------------------------------------------------------------------
\section*{Part B: Experimental Design (45 points)}

\begin{question}[Testing dynamic complementarity in the field]
Design a field experiment in a low-income setting that tests whether
early and later child investments are complements --- not merely whether
each program "works." Specify:
\begin{enumerate}[label=(\alph*)]
  \item \textbf{(6 pts)} Setting, subject pool, and recruitment: the
        two investment programs (an early-childhood program, ages 1--3,
        and a later program, ages 5--7) and why your setting makes the
        complementarity question policy-relevant.
  \item \textbf{(12 pts)} Treatments: the cross-randomized 2$\times$2
        (neither / early only / late only / both), what each comparison
        identifies, and why the \emph{interaction} is the estimand ---
        including the prediction each technology (complements vs.\
        substitutes) makes for the ``both'' cell relative to the sum of
        the single-program effects.
  \item \textbf{(9 pts)} Measurement: child skill assessments at
        multiple ages AND parental investment responses (time, money,
        attention) --- explain why measuring parental responses is not
        optional in this design (Part A(1)(d) is the reason), and how
        you will distinguish reinforcement from compensation in the
        data.
  \item \textbf{(9 pts)} Hypotheses and power: state the interaction
        hypothesis under each technology, and --- using Module 6's
        lesson that interactions need roughly 4$\times$ the sample of a
        main effect --- give per-cell and total sample sizes anchored to
        a stated main-effect size (e.g.\ 0.25 SD per program, typical
        of early-childhood interventions; any reasoned anchor
        accepted). Note the follow-up horizon problem and your interim
        strategy.
  \item \textbf{(9 pts)} Ethics with a vulnerable population --- this
        is graded as designed substance, not boilerplate: equipoise for
        each cell (why is ``neither'' ethical here?), what the control
        group receives, consent under low literacy, local approval and
        benefit sharing, and one design change you would accept at an
        ethics board's request even at a power cost.
\end{enumerate}
\end{question}

% ---------------------------------------------------------------------------
\section*{Part C: Silicon Pilot Interpretation $\langle$AI$\rangle$
(25 points)}

\begin{question}[The population that matched the mean]
You run \texttt{module7-simulation-lab.ipynb}. Your synthetic
rural-Kenya population's mean dictator-game giving lands within one
percentage point of the published human estimate you entered as the
calibration target (you did not tune the population to the target; the
match arose on the first run). Your lab partner declares the population
``calibrated'' and proposes piloting a cash-transfer instrument on it to
forecast the transfer's consumption response for the grant proposal.

\begin{enumerate}[label=(\alph*)]
  \item \textbf{(8 pts)} Using the calibration hierarchy (marginals →
        behavioral moment → gradients/heterogeneity → treatment
        responses), explain what the matched mean does and does not
        establish, and give two LLM-specific mechanisms that can
        produce a matched mean WITHOUT the population being calibrated
        in any useful sense.
  \item \textbf{(9 pts)} Design the diagnostic battery for the next two
        rungs: (i) a gradient check (which demographic gradients, from
        which human literature, and what pattern confirms vs.\ refutes
        calibration) and (ii) an out-of-task check (a second game or
        elicitation with published cross-cultural data). State what
        result pattern would justify which uses of the population. Run
        it if API budget allows and report; otherwise state predictions
        precisely.
  \item \textbf{(8 pts)} The grant proposal: state one thing the
        synthetic population CAN legitimately contribute to the
        cash-transfer study (and why), one thing it CANNOT (and the
        circularity that makes rung 4 unreachable from inside the
        sim), and the sentence you would write in the proposal's
        methods section describing the silicon pilot's role honestly.
\end{enumerate}
Attach your AI transcripts and, if you ran code, the results CSV.
\end{question}

\end{document}
